
\documentclass[a4paper,11pt]{article}
\usepackage{fontspec}
\setmainfont{Noto Sans}
\usepackage[french]{babel}
\usepackage{microtype}
\usepackage[top=9mm,bottom=8mm,left=11mm,right=11mm]{geometry}
\usepackage{xcolor}
\usepackage{tikz}
\usetikzlibrary{arrows.meta,calc,positioning}
\usepackage[most]{tcolorbox}
\usepackage{ulem}
\pagestyle{empty}

\definecolor{Blue}{HTML}{356FA8}
\definecolor{BlueDark}{HTML}{244D78}
\definecolor{BlueLight}{HTML}{EAF2F9}
\definecolor{BlueBorder}{HTML}{B9D3E3}
\definecolor{Violet}{HTML}{7A5AA6}
\definecolor{VioletDark}{HTML}{563A7A}
\definecolor{VioletLight}{HTML}{F2ECF8}
\definecolor{VioletBorder}{HTML}{D3C4E2}
\definecolor{Neutral}{HTML}{6C7378}
\definecolor{NeutralDark}{HTML}{42494D}
\definecolor{NeutralLight}{HTML}{F0F2F3}
\definecolor{NeutralBorder}{HTML}{C9CED1}
\definecolor{Ink}{HTML}{172A32}
\definecolor{Grey}{HTML}{64757D}
\definecolor{Red}{HTML}{D52F35}
\definecolor{RedLight}{HTML}{FCEBED}

\tcbset{
 enhanced,boxrule=.6pt,arc=2.8mm,
 left=3.7mm,right=3.7mm,top=1.75mm,bottom=1.75mm,
 boxsep=1mm,coltext=Ink
}
\newtcolorbox{bluebox}[1][]{colback=BlueLight,colframe=BlueBorder,#1}
\newtcolorbox{violetbox}[1][]{colback=VioletLight,colframe=VioletBorder,#1}
\newtcolorbox{neutralbox}[1][]{colback=NeutralLight,colframe=NeutralBorder,#1}
\newtcolorbox{warnbox}[1][]{colback=RedLight,colframe=Red,
 left=2.7mm,right=2.7mm,top=1.25mm,bottom=1.25mm,#1}

\begin{document}
\begin{tikzpicture}[remember picture,overlay]
\draw[draw=NeutralBorder,line width=.9pt,rounded corners=1.7mm]
([xshift=1mm,yshift=-1mm]current page.north west)--([xshift=-1mm,yshift=-1mm]current page.north east)--
([xshift=-1mm,yshift=1mm]current page.south east)--([xshift=1mm,yshift=1mm]current page.south west)--cycle;
\fill[Neutral]([xshift=1mm,yshift=-1mm]current page.north west) rectangle ([xshift=5.2mm,yshift=1mm]current page.south west);
\end{tikzpicture}
\vspace*{1mm}\hspace*{6mm}{\color{NeutralDark}\fontsize{25}{27}\selectfont\bfseries INFINITIU E PARTICIPI PRESENT}
\par\vspace{.5mm}\hspace*{6mm}{\color{Grey}\large L'infinitif et le participe présent}
\par\vspace{2.5mm}\noindent\begin{minipage}{.94\textwidth}
\begin{neutralbox}\textbf{\color{NeutralDark}FORMAS}\\[1mm]\centering
{\Large\textbf{INFINITIU \quad / \quad PARTICIPI PRESENT}}\\[1mm]
\small \textit{cantar} \quad / \quad \textit{cantant}
\end{neutralbox}
\vspace{2.5mm}\begin{center}\begin{tikzpicture}[>=Latex]
\draw[Ink!65,line width=.8pt](0,0)--(15.1,0);
\node[below=2mm]at(3.2,0){\small NOM D'ACCION};\node[below=2mm]at(11.7,0){\small ACCION EN CORS};
\draw[Neutral,line width=2.5pt,->](2.4,0)--(5.3,0);
\draw[Neutral,line width=2.5pt,->](9.8,0)--(13.5,0);
\node[above=3mm,align=center,text width=5.8cm]at(4,0){\small\bfseries infinitiu : nominar l'accion};
\node[above=3mm,align=center,text width=5.8cm]at(11.7,0){\small\bfseries participi present : caracterizar l'accion};
\end{tikzpicture}\end{center}
\vspace{1.5mm}
\begin{neutralbox}\textbf{\color{NeutralDark}INFINITIU}\\[1mm]
L'infinitiu es la forma de basa del vèrb : \textbf{cantar, finir, vendre...}\\
Pòt èsser emplegat après d'autres vèrbs e dins de construccions ont l'accion es considerada coma una nocion o una activitat.
\end{neutralbox}
\vspace{2mm}
\begin{neutralbox}\textbf{\color{NeutralDark}PARTICIPI PRESENT}\\[1mm]
Lo participi present es una forma en \textbf{-ant} : \textit{cantant, parlant, escrivent...}\\
Servís notamment a caracterizar una accion o un èstre e pòt participar a de construccions adjectivalas o circonstancialas.
\end{neutralbox}
\vspace{2mm}
\begin{minipage}{.485\textwidth}\begin{neutralbox}[height=3.15cm]
\textbf{\color{NeutralDark}EXEMPLES}\\[1mm]
\textit{Ai enveja de cantar.}\\
\textit{Vesi un òme corrent.}\\
\textit{Parlant aital, s'es explicat.}
\end{neutralbox}\end{minipage}\hfill
\begin{minipage}{.485\textwidth}\begin{warnbox}[height=3.15cm]
\textbf{\color{Red}ATENCION}\\[1mm]
Lo participi present es pas l'equivalent sistematic del \textit{-ing} anglés : son usatge e sas construccions dependon de la sintaxi occitana.
\end{warnbox}\end{minipage}
\vspace{2mm}\begin{neutralbox}\textbf{\color{NeutralDark}A RETÉNER}\\[1mm]
L'\textbf{infinitiu} nomina l'accion ; lo \textbf{participi present} permet de la caracterizar o de l'integrar dins una construccion.
\end{neutralbox}
\vspace{1.5mm}\centering{\small\color{Grey}\textbf{Exemples de formes :} parlar, finir, vendre / parlant, finissent, vendent...}
\vspace{1mm}{\scriptsize\color{Grey}Referéncia : occitan languedocian, graphia classica, vèrb'Òc / Lo Congrès.}
\end{minipage}\end{document}
